\section{Discussion}
\label{sec:discussion}

\noindent In this study, we developed an MR framework that uses population-scale single-cell multiomics data to estimate the causal effects of chromatin accessibility on gene expression. Rather than relying on marginal associations across cells or individuals, the framework uses cell-type-specific \textit{cis}-caQTL and \textit{cis}-eQTL summary statistics as genetic instruments. Applying it to three immune cell types in CIMA, we compared the MR estimates with two marginal measures of peak--gene association. The two approaches agreed for many links but diverged for others: the MR estimates were largely unrelated to correlations across metacells, which pool several immune cell types, and agreed more closely, though still only modestly, with correlations across donors within a cell type. These patterns illustrate that correlation does not always reflect causation.

\noindent We assessed the regulatory relevance of MR-supported links using resources independent of the QTL summary statistics: enhancer annotations, blood eQTLs, promoter capture Hi-C, H3K27ac HiChIP contacts, and the eRegulons inferred by CIMA. Links supported by MR were enriched for every external reference, generally with larger fold enrichment than links supported by marginal correlations, and showed higher odds of eRegulon membership. Notably, the enrichment for external references extended to MR-only links, which were detected by MR but by neither marginal method; genes involved in these links were enriched for transcriptional regulation, RNA processing, and protein ubiquitination. Finally, peaks in MR-only links showed nominal heritability enrichment for several immune-mediated diseases, and individual MR-only links nominated candidate target genes and cellular contexts for disease-associated variants near \textit{ERAP1}, \textit{ERBB3}, and \textit{THBS3}, providing hypotheses for experimental follow-up.

\noindent Our study is limited in several ways. First, the \textit{cis}-caQTL and \textit{cis}-eQTL summary statistics were derived from the same CIMA cohort, whereas two-sample MR typically uses SNP--exposure and SNP--outcome associations from independent cohorts so that their estimation errors are uncorrelated \cite{burgess2016overlap}. This sample overlap is closely tied to the No Measurement Error (NOME) assumption of IVW and MR-Egger, which treats the estimation error in the SNP--exposure associations as negligible. Both problems act through the same quantity, the covariance between the estimated SNP--exposure associations and the residuals of the MR regression (\nameref{S1_Text}, Section~4). Without sample overlap, measurement error in the SNP--exposure associations attenuates causal estimates toward the null; with complete overlap, the bias is directed toward the marginal association. In either case, sample overlap introduces bias only when NOME is violated, and the bias is approximately inversely proportional to instrument strength, as measured by the $F$-statistic. We therefore used only strong instruments (nominal \textit{cis}-caQTL $P<10^{-8}$), whose $F$-statistics far exceed the conventional cutoff of 10; this keeps the bias small for IVW, although MR-Egger is more sensitive to measurement error \cite{bowden2016nome}. In addition, a sensitivity analysis using nonsignificant (null) SNPs showed that their caQTL and eQTL $z$-scores were only weakly correlated (\nameref{S1_Fig}), indicating that the error correlation induced by sample overlap is small on average. This average is dominated by the many pairs without a causal effect; for pairs with a nonzero causal effect, the residuals of accessibility and expression are correlated by construction, and protection against bias comes mainly from instrument strength. The stringent instrument threshold may exclude weaker genetic signals and reduce statistical power. Weak-instrument-robust estimators, such as limited information maximum likelihood (LIML) and Fuller estimators, typically require individual-level genotype data and rely on the strict exclusion restriction that variants affect gene expression solely through the target peak's accessibility, an assumption rarely met in \textit{cis}-regulatory mapping. As population-scale single-cell multiomics and QTL resources continue to expand, future studies can draw caQTL and eQTL summary statistics from independent cohorts to avoid sample overlap altogether.
